\documentclass[a4paper,11pt]{article}

% ============================================================
%  Le Collé — Algèbre linéaire (compléments de deuxième année)
%  © 2026 Ahmed Amine Ouqlani
% ============================================================

% ---------- Polices ----------
\usepackage[utf8]{inputenc}
\usepackage[T1]{fontenc}
\usepackage[sc]{mathpazo}          % texte et mathématiques : Palatino
\usepackage[scaled=0.88]{helvet}   % titres et repères : linéale
\linespread{1.07}
\usepackage{amsmath,amssymb,amsthm,mathrsfs,stmaryrd}
\usepackage{microtype}

% ---------- Mise en page ----------
\usepackage[a4paper,top=2.3cm,bottom=2.3cm,left=2.4cm,right=2.4cm,
            headheight=14pt,headsep=0.8cm,footskip=1.1cm]{geometry}
\setlength{\parindent}{0pt}
\setlength{\parskip}{0.45em plus 0.15em minus 0.1em}
\raggedbottom
\emergencystretch=1.5em
\clubpenalty=10000 \widowpenalty=10000 \displaywidowpenalty=10000 \brokenpenalty=10000

% ---------- Couleurs : un bleu de nuit, un ocre, des gris ----------
\usepackage{xcolor}
\definecolor{Ink}{HTML}{1C2533}
\definecolor{Bleu}{HTML}{1F3A5F}
\definecolor{BleuMoyen}{HTML}{4A6A8F}
\definecolor{BleuClair}{HTML}{9DB2C9}
\definecolor{Fond}{HTML}{F2F4F7}
\definecolor{Ocre}{HTML}{B07A26}
\definecolor{Brique}{HTML}{A23B2A}
\definecolor{Gris}{HTML}{5E6672}
\definecolor{Filet}{HTML}{C5CCD6}

% ---------- Packages ----------
\usepackage{enumitem}
\setlist{itemsep=0.25em,topsep=0.35em,parsep=0pt}
\setlist[itemize]{label=\textcolor{BleuMoyen}{\small\textbullet},leftmargin=1.4em}
\setlist[enumerate]{leftmargin=1.8em,label=\textsf{\small\color{Bleu}\arabic*.}}
\usepackage{booktabs,array,multicol}
\usepackage{needspace,float}
\usepackage[font=small,labelfont={sf,color=Bleu},labelsep=period,skip=6pt]{caption}
\usepackage{tikz}
\usetikzlibrary{arrows.meta,calc,decorations.markings}
\usepackage{pgfplots}
\pgfplotsset{compat=1.17}
\usepackage{tkz-tab}
\usepackage[most]{tcolorbox}
\usepackage{titlesec,titletoc}
\usepackage{fancyhdr,lastpage,etoolbox}
\usepackage[colorlinks=true,linkcolor=Bleu,urlcolor=Bleu,citecolor=Bleu,
            pdftitle={Le Collé — Algèbre linéaire},pdfauthor={Ahmed Amine Ouqlani}]{hyperref}

\numberwithin{equation}{section}
\numberwithin{figure}{section}
\numberwithin{table}{section}

% ---------- Titres ----------
% numéro du chapitre dans le parcours PT (affiché), distinct de la numérotation interne
\newcommand{\grandnum}{\thesection}\newcommand{\ongletcouleur}{Bleu}
% Chapitres (sections) : grand numéro, titre en linéale, filet ocre
\titleformat{\section}[display]
  {\normalfont\sffamily\color{Ink}}
  {{\small\color{Gris}\textls[200]{\MakeUppercase{Chapitre}}}\\[-0.2em]%
   {\fontsize{64}{64}\selectfont\color{BleuClair}\grandnum}}
  {0.4em}
  {\fontsize{25}{30}\selectfont\bfseries}
  [\vspace{0.5em}{\color{Ocre}\rule{3.2cm}{1.4pt}}]
\titlespacing*{\section}{0pt}{0pt}{1.6em}
\titleformat{name=\section,numberless}[display]
  {\normalfont\sffamily\color{Ink}}{}{0pt}{\fontsize{25}{30}\selectfont\bfseries}
  [\vspace{0.5em}{\color{Ocre}\rule{3.2cm}{1.4pt}}]
\titlespacing*{name=\section,numberless}{0pt}{0pt}{1.6em}

\titleformat{\subsection}
  {\needspace{5\baselineskip}\normalfont\sffamily\Large\bfseries\color{Ink}}
  {\color{Bleu}\thesubsection}{0.8em}{}
\titlespacing*{\subsection}{0pt}{1.8em}{0.5em}
\titleformat{\subsubsection}
  {\needspace{5\baselineskip}\normalfont\sffamily\large\bfseries\color{Ink}}
  {\color{Bleu}\thesubsubsection}{0.7em}{}
\titlespacing*{\subsubsection}{0pt}{1.6em}{0.4em}
\titleformat{\paragraph}[runin]{\normalfont\sffamily\bfseries\color{Ink}}{}{0pt}{}[.]
\titlespacing*{\paragraph}{0pt}{0.9em}{0.6em}
\setcounter{secnumdepth}{3}
\setcounter{tocdepth}{2}

% ---------- Table des matières ----------
\titlecontents{section}[0pt]{\addvspace{0.7em}\sffamily\bfseries\color{Ink}}
  {}{}{\hfill\contentspage}[\addvspace{0.2em}]
\titlecontents{subsection}[2.2em]{\small\linespread{1}\selectfont}
  {\makebox[2.6em][l]{\color{Gris}\thecontentslabel}}{}{\titlerule*[0.7pc]{.}\contentspage}

% ---------- En-têtes et pieds de page ----------
\pagestyle{fancy}
\fancyhf{}
\renewcommand{\sectionmark}[1]{\markboth{#1}{}}
\fancyhead[L]{\footnotesize\sffamily\color{Gris}\leftmark}
\fancyhead[R]{\footnotesize\sffamily\scshape\color{Gris}Le Collé}
\fancyfoot[C]{\small\sffamily\color{Ink}\thepage}
\fancyfoot[L]{\scriptsize\sffamily\hyperlink{toc}{\color{Filet}sommaire}}
\renewcommand{\headrulewidth}{0.4pt}
\renewcommand{\headrule}{{\color{Filet}\hrule height\headrulewidth width\headwidth}}
\renewcommand{\footrulewidth}{0pt}
\fancypagestyle{plain}{\fancyhf{}\fancyfoot[C]{\small\sffamily\color{Ink}\thepage}\renewcommand{\headrulewidth}{0pt}}

% ---------- Énoncés ----------
\newtheoremstyle{enonce}{0pt}{0pt}{}{}{\sffamily\bfseries\color{Bleu}}{.}{0.5em}
  {\thmname{#1}\thmnumber{~#2}\thmnote{\ \normalfont\sffamily\color{Ink}(#3)}}
\newtheoremstyle{exemple}{0.9em}{0.9em}{}{}{\sffamily\bfseries\color{BleuMoyen}}{.}{0.5em}
  {\thmname{#1}\thmnumber{~#2}\thmnote{\ \normalfont\sffamily\color{Ink}(#3)}}
\newtheoremstyle{remarque}{0.7em}{0.7em}{}{}{\sffamily\itshape}{.}{0.5em}{\thmname{#1}\thmnote{\ (#3)}}
\theoremstyle{enonce}
\newtheorem{definition}{Définition}[section]
\newtheorem{theoreme}[definition]{Théorème}
\newtheorem{proposition}[definition]{Proposition}
\newtheorem{corollaire}[definition]{Corollaire}
\theoremstyle{exemple}
\newtheorem{exemple}[definition]{Exemple}
\theoremstyle{remarque}
\newtheorem*{remarque}{Remarque}

\tcbset{
  defi/.style={enhanced,breakable,colback=white,colframe=white,boxrule=0pt,arc=0pt,
     borderline west={1.2pt}{0pt}{Bleu},left=11pt,right=4pt,top=3pt,bottom=3pt,
     before skip=1.1em,after skip=1.1em},
  thm/.style={enhanced,breakable,colback=Fond,colframe=Fond,boxrule=0pt,arc=0pt,
     borderline west={2.4pt}{0pt}{Bleu},left=11pt,right=8pt,top=6pt,bottom=6pt,
     before skip=1.2em,after skip=1.2em}}
\tcolorboxenvironment{definition}{defi}
\tcolorboxenvironment{theoreme}{thm}
\tcolorboxenvironment{proposition}{thm}
\tcolorboxenvironment{corollaire}{thm}

\renewcommand{\proofname}{Démonstration}
\renewcommand{\qedsymbol}{\textcolor{BleuClair}{$\blacksquare$}}
\makeatletter
\renewenvironment{proof}[1][\proofname]{\par\pushQED{\qed}\normalfont
  \topsep0.5em\relax\trivlist\item[\hskip\labelsep\sffamily\itshape\color{Gris}#1.]\ignorespaces}
  {\popQED\endtrivlist\@endpefalse}
\makeatother

% ---------- Méthode : filets haut et bas, titre en bandeau ----------
\makeatletter\newcommand{\methlabelfix}{\protected@edef\@currentlabel{\thetcbcounter}}\makeatother
\newtcolorbox[auto counter,number within=section]{methode}[2][]{%
  enhanced,unbreakable,colback=white,colframe=white,boxrule=0pt,arc=0pt,
  borderline north={0.9pt}{0pt}{Bleu},borderline south={0.9pt}{0pt}{Bleu},
  left=4pt,right=4pt,top=8pt,bottom=8pt,before skip=1.4em,after skip=1.4em,
  before upper={\methlabelfix\phantomsection\addcontentsline{met}{metentry}{\protect\makebox[2.6em][l]{\thetcbcounter}#2}%
     {\sffamily\bfseries\small\color{white}\colorbox{Bleu}{\strut\ MÉTHODE~\thetcbcounter\ }}%
     \ \ {\sffamily\bfseries\color{Bleu}#2}\par\smallskip},#1}

% ---------- Repères pédagogiques : un filet de couleur, un mot-clé ----------
\newcommand{\definirbloc}[3]{% nom, couleur, libellé
  \newtcolorbox{#1}[1][]{enhanced,breakable,blanker,
    borderline west={1.4pt}{0pt}{#2},left=10pt,right=2pt,top=1pt,bottom=1pt,
    before skip=1em,after skip=1em,
    before upper={{\sffamily\bfseries\footnotesize\color{#2}\textls[120]{\MakeUppercase{#3}}}\quad},##1}}
\definirbloc{reflexe}{Bleu}{Réflexe}
\definirbloc{astuce}{Ocre}{Astuce}
\definirbloc{attention}{Brique}{Attention}
\definirbloc{enpratique}{BleuMoyen}{En pratique}
\definirbloc{plusloin}{Gris}{Pour aller plus loin}
\definirbloc{strategie}{Bleu}{Stratégie}
\definirbloc{lecture}{BleuMoyen}{Lire le schéma}
% compatibilité avec les parties reprises
\newenvironment{garde}{\begin{attention}}{\end{attention}}
\newenvironment{complementHP}{\begin{plusloin}}{\end{plusloin}}
\newenvironment{interpretation}{\begin{enpratique}}{\end{enpratique}}

% ---------- Exercices, corrigés ----------
\newcounter{exo}[section]
\renewcommand{\theexo}{\thesection.\arabic{exo}}
\newcommand{\niveau}[1]{{\color{Ocre}\footnotesize
  \ifcase#1\or$\bullet\circ\circ\circ\circ$\or$\bullet\bullet\circ\circ\circ$\or$\bullet\bullet\bullet\circ\circ$%
  \or$\bullet\bullet\bullet\bullet\circ$\or$\bullet\bullet\bullet\bullet\bullet$\fi}}
\newcommand{\serie}[1]{\par\needspace{8\baselineskip}\vspace{1.5em}%
  \noindent{\sffamily\bfseries\color{Bleu}#1}\par\nobreak\vspace{-0.4em}%
  \noindent{\color{Filet}\rule{\linewidth}{0.4pt}}\par\nobreak\vspace{0.1em}}
\newenvironment{exo}[3]{%
  \par\needspace{5\baselineskip}\vspace{0.75em}\refstepcounter{exo}\label{e:#1}%
  \def\exolab{#1}%
  \noindent{\sffamily\bfseries\color{Bleu}Exercice~\theexo}\quad\niveau{#2}%
  \hfill{\footnotesize\sffamily\color{Gris}#3}\par\nobreak\vspace{0.15em}}%
  {\par\nopagebreak\noindent\hfill{\scriptsize\sffamily\hyperref[c:\exolab]{\color{Gris}corrigé p.~\pageref{c:\exolab}}}\par}
\newenvironment{corrige}[1]{%
  \par\needspace{6\baselineskip}\vspace{1em}\phantomsection\label{c:#1}%
  \noindent{\sffamily\bfseries\color{Bleu}Exercice~\ref{e:#1}}\quad
  {\scriptsize\sffamily\hyperref[e:#1]{\color{Gris}énoncé p.~\pageref{e:#1}}}\par\nobreak\vspace{0.2em}}%
  {\hfill\textcolor{BleuClair}{$\square$}\par}
\newcommand{\etape}[1]{\par\smallskip\noindent{\sffamily\bfseries\small\color{BleuMoyen}#1.}\ }
\newenvironment{solution}{\par\noindent{\sffamily\bfseries\small\color{BleuMoyen}Solution.}\ }{\hfill\textcolor{BleuClair}{$\square$}\par}
\newcommand{\chercher}{\par\smallskip\noindent{\color{Filet}\rule[0.5ex]{1.2cm}{0.4pt}}\
  {\footnotesize\sffamily\itshape\color{Gris}laisser chercher l'étudiant avant de poursuivre}\par\smallskip}

% ---------- Guide de séance ----------

\newcommand{\marge}[1]{\leavevmode\marginpar{\raggedleft\scriptsize\sffamily\color{Gris}\textls[100]{\MakeUppercase{#1}}}\ignorespaces}
\newcommand{\Etape}[4]{%
  \subsection[Étape #1 --- #3 {\color{Gris}(#2)}]{#3}\label{etape:#1}%
  \noindent{\small\sffamily\color{Gris}{\color{Ocre}\bfseries Étape #1}\quad·\quad #2\quad·\quad #4}\par\medskip}
\newcommand{\question}[1]{\par\smallskip\noindent{\sffamily\bfseries\small\color{BleuMoyen}Question.}\ \textit{#1}\par\smallskip}
\newcommand{\rappel}[1]{\par\smallskip\noindent\phantomsection\addcontentsline{rap}{rapentry}{#1}%
  {\sffamily\bfseries\scriptsize\colorbox{Fond}{\color{Gris}\strut RAPPEL}}\ {\sffamily\small\color{Gris}#1.}\ }
\newcommand{\pourquoi}{\par\noindent{\sffamily\bfseries\small\color{Ocre}Pourquoi ?}\ }

% Ouverture de chapitre : \ouverture{objectifs}{prérequis}{durée}
\newcommand{\ouverture}[3]{%
  \begin{tcolorbox}[enhanced,colback=Fond,colframe=Fond,arc=0pt,boxrule=0pt,left=12pt,right=12pt,top=10pt,bottom=10pt,
     before skip=1.2em,after skip=1.6em]
  \begin{minipage}[t]{0.56\linewidth}
  {\sffamily\bfseries\small\color{Bleu}OBJECTIFS}\par\vspace{0.3em}\small #1
  \end{minipage}\hfill
  \begin{minipage}[t]{0.38\linewidth}
  {\sffamily\bfseries\small\color{Bleu}PRÉREQUIS}\par\vspace{0.3em}\small #2
  \ifx\relax#3\relax\else\par\vspace{0.6em}{\sffamily\bfseries\small\color{Bleu}DURÉE}\par\vspace{0.2em}\small #3\fi
  \end{minipage}
  \end{tcolorbox}}
% Fin de chapitre
\newenvironment{essentiel}{\par\needspace{12\baselineskip}\vspace{1.5em}%
  \begin{tcolorbox}[enhanced,unbreakable,colback=white,colframe=white,boxrule=0pt,arc=0pt,
   borderline north={1.4pt}{0pt}{Ocre},left=2pt,right=2pt,top=8pt,bottom=4pt]
  {\sffamily\bfseries\color{Ocre}L'ESSENTIEL DU CHAPITRE}\par\medskip}{\end{tcolorbox}}

% ---------- Notations ----------
\newcommand{\K}{\mathbb{K}}
\newcommand{\R}{\mathbb{R}}
\newcommand{\C}{\mathbb{C}}
\newcommand{\N}{\mathbb{N}}
\newcommand{\Z}{\mathbb{Z}}
\newcommand{\Mn}[1]{\mathcal{M}_{n}(#1)}
\newcommand{\GLn}[1]{\mathrm{GL}_n(#1)}
\newcommand{\LE}{\mathcal{L}(E)}
\newcommand{\rg}{\operatorname{rg}}
\newcommand{\vect}{\operatorname{Vect}}
\newcommand{\Mat}{\operatorname{Mat}}
\newcommand{\Id}{\operatorname{Id}}
\newcommand{\sh}{\operatorname{sh}}
\newcommand{\ch}{\operatorname{ch}}
\renewcommand{\th}{\operatorname{th}}
\DeclareMathOperator{\argsh}{argsh}
\newcommand{\vv}[1]{\overrightarrow{#1}}
\newcommand{\ve}[1]{\vec{#1}}
\newcommand{\ient}[2]{\llbracket #1,#2\rrbracket}
\newcommand{\dt}[1]{\left|\begin{matrix}#1\end{matrix}\right|}
\newcommand{\dis}{\displaystyle}

% ---------- Figures ----------
\pgfplotsset{trig format plots=rad}
\pgfplotsset{courbe/.style={axis lines=middle,axis equal image,
  axis line style={Gris,-{Stealth[length=5pt]}},
  xlabel={$x$},ylabel={$y$},ticklabel style={font=\scriptsize\sffamily,color=Gris},
  every axis x label/.style={at={(ticklabel* cs:1)},anchor=west,font=\small},
  every axis y label/.style={at={(ticklabel* cs:1)},anchor=south,font=\small},
  samples=300,smooth,width=0.6\linewidth}}
\tikzset{arcetude/.style={Bleu,line width=1.25pt},
         reste/.style={BleuClair,line width=1pt},
         tgte/.style={Ocre,line width=0.9pt},
         asy/.style={Gris,dashed,line width=0.6pt},
         pt/.style={circle,fill=Ink,inner sep=1.2pt}}

% ---------- Navigation ----------
\makeatletter
\newcommand*\l@rapentry{\@dottedtocline{1}{0em}{0em}}
\newcommand*\l@metentry{\@dottedtocline{1}{0em}{0em}}
\newcommand{\listemethodes}{\@starttoc{met}}
\newcommand{\listerappels}{\@starttoc{rap}}
\makeatother
% onglets de chapitre sur le bord droit
\usepackage{eso-pic}
\newif\ifonglets \ongletsfalse
\newcommand{\dessineronglet}{%
  \edef\ongcol{\ongletcouleur}%
  \edef\ongnum{\grandnum}%
  \pgfmathsetmacro{\ongy}{3.2+2.3*(\ongnum-1)}%
  \begin{tikzpicture}[remember picture,overlay]
  \fill[\ongcol] ([xshift=-0.75cm,yshift=-\ongy cm]current page.north east)
     rectangle ([yshift=-\ongy cm-1.9cm]current page.north east);
  \node[text=white,font=\sffamily\bfseries\large] at ([xshift=-0.37cm,yshift=-\ongy cm-0.95cm]current page.north east)
     {\ongnum};
  \end{tikzpicture}}
\AddToShipoutPictureBG{\ifonglets\ifnum\value{section}>0 \dessineronglet\fi\fi}
% blocs de navigation et d'enseignement
\definirbloc{commencer}{Gris}{Si vous commencez ici}
\definirbloc{horsprogramme}{Brique}{Hors programme}
\newcommand{\HP}{\,{\sffamily\bfseries\scriptsize\textcolor{Brique}{[HP]}}}
\newcommand{\source}[1]{{\footnotesize\sffamily\color{Gris}\,(#1)}}
\newcommand{\coupdepouce}[1]{\par\vspace{-0.2em}\noindent{\footnotesize\sffamily\color{Gris}%
  \textbf{Coup de pouce} (à donner si l'étudiant bloque) : #1}\par}
\newcommand{\pointarret}[1]{\par\vspace{1em}\noindent{\color{Filet}\rule[0.5ex]{\linewidth}{0.4pt}}\par\nobreak
  \noindent{\footnotesize\sffamily\color{Ocre}\bfseries POINT D'ARRÊT POSSIBLE}\quad{\footnotesize\sffamily\color{Gris}#1}\par\vspace{0.6em}}
\newcommand{\demarche}[1]{\noindent{\sffamily\bfseries\small\color{BleuMoyen}Démarche.}\ \textit{#1}\par\smallskip}

% ---------- Corrigés : idées, astuces, pièges, à retenir ----------
\newcommand{\choix}[1]{\par\noindent{\sffamily\bfseries\small\color{BleuMoyen}Choix de la méthode.}\ #1\par\smallskip}
\newcommand{\gain}[1]{\par\noindent{\sffamily\bfseries\small\color{Ocre}Astuce.}\ #1\par}
\newcommand{\piege}[1]{\par\noindent{\sffamily\bfseries\small\color{Brique}Piège.}\ #1\par}
\newtcolorbox{retenirbox}{enhanced,breakable,colback=Fond,colframe=Fond,arc=0pt,boxrule=0pt,
  borderline west={2pt}{0pt}{Ocre},left=8pt,right=6pt,top=4pt,bottom=4pt,before skip=0.7em,after skip=0.5em}
\newcommand{\retenir}[1]{\begin{retenirbox}{\sffamily\bfseries\footnotesize\color{Ocre}À RETENIR}\quad #1\end{retenirbox}}
\newcommand{\TD}[1]{{\color{Ocre}\bfseries TD\,#1}\ ·\ }
\newcommand{\q}[1]{\par\smallskip\noindent{\sffamily\bfseries\color{Bleu}#1}\ }
\providecommand{\iddots}{\mathinner{\mkern1mu\raise1pt\hbox{.}\mkern2mu\raise4pt\hbox{.}\mkern2mu\raise7pt\hbox{.}\mkern1mu}}


\begin{document}
\renewcommand{\grandnum}{1}
\pagenumbering{roman}
\begin{titlepage}
\thispagestyle{empty}
\begin{tikzpicture}[remember picture,overlay]
  \fill[Bleu] (current page.north west) rectangle ([yshift=-15.5cm]current page.north east);
  \node[anchor=north west,text=white,font=\sffamily\small] at ([xshift=2.6cm,yshift=-3.2cm]current page.north west)
     {\textls[250]{MATHÉMATIQUES\ \ ·\ \ CPGE PT\ \ ·\ \ DEUXIÈME ANNÉE}};
  \node[anchor=north west,text=white] at ([xshift=2.5cm,yshift=-4.4cm]current page.north west)
     {\fontsize{62}{66}\selectfont\scshape Le Collé};
  \draw[Ocre,line width=1.6pt] ([xshift=2.6cm,yshift=-7.3cm]current page.north west) -- ++(4.2cm,0);
  \node[anchor=north west,text=white,font=\sffamily\Large\bfseries,align=left]
     at ([xshift=2.6cm,yshift=-8.0cm]current page.north west) {Algèbre linéaire};
  \node[anchor=north west,text=white!90!Bleu,font=\sffamily\large,align=left]
     at ([xshift=2.6cm,yshift=-9.0cm]current page.north west) {Sommes directes, hyperplans, trace\\[0.15cm]Sous-espaces stables et matrices par blocs};
  \node[anchor=north west,text=Ink,font=\sffamily\small,align=left,text width=12cm]
     at ([xshift=2.6cm,yshift=-17.4cm]current page.north west)
     {\textcolor{Bleu}{\bfseries Cours, méthodes et exercices corrigés}\\[0.3cm]\textcolor{Gris}{Les compléments d'algèbre linéaire de deuxième année : le langage des chapitres \emph{Déterminants} et \emph{Réduction}.}};
  \node[anchor=south west,text=Ink,font=\sffamily\large] at ([xshift=2.6cm,yshift=3.0cm]current page.south west) {Ahmed Amine Ouqlani};
  \draw[Filet] ([xshift=2.6cm,yshift=2.6cm]current page.south west) -- ([xshift=-2.4cm,yshift=2.6cm]current page.south east);
\end{tikzpicture}
\end{titlepage}
\thispagestyle{plain}
\section*{Avant-propos}
\markboth{Avant-propos}{}
\providecommand{\Ker}{\operatorname{Ker}}\renewcommand{\Im}{\operatorname{Im}}\providecommand{\tr}{\operatorname{tr}}\providecommand{\Sp}{\operatorname{Sp}}
Ce livret ouvre le parcours d'algèbre de deuxième année. Il rassemble les \emph{compléments d'algèbre linéaire} du programme de
PT : produit et somme de sous-espaces, somme directe, hyperplans, trace, sous-espaces stables et matrices par blocs. Ce sont les
outils dont ont besoin les deux chapitres suivants : \emph{Déterminants} (chapitre~2) et surtout \emph{Réduction} (chapitre~3),
qui renvoie ici au lieu de les redémontrer.

La première partie rappelle, sans démonstration, les résultats de première année utilisés ensuite (rang, théorème du rang,
matrice d'une application linéaire, changement de base) : on la parcourt pour vérifier qu'ils sont acquis. Chaque résultat
nouveau est énoncé avec ses hypothèses, illustré par un exemple, et chaque partie se termine par des exercices, corrigés à la fin
du livret. La présentation, la fiche, le TD et la sélection de problèmes du chapitre suivent la même numérotation (§1.1 à §1.6).

\vfill
{\footnotesize\sffamily\color{Gris}\textcopyright~2026 Ahmed Amine Ouqlani. Tous droits réservés.\par}
\clearpage
\thispagestyle{plain}
\hypertarget{toc}{}
\renewcommand{\contentsname}{Sommaire}
\tableofcontents
\clearpage
\pagenumbering{arabic}
\ongletstrue
% =====================================================================
%                CHAPITRE — ALGÈBRE LINÉAIRE : COMPLÉMENTS
% =====================================================================
\clearpage
\section{Algèbre linéaire}\label{sec:al}

En première année, on a appris à représenter une application linéaire par une matrice. La deuxième année pose une question plus
fine : comment \emph{découper} un espace en morceaux simples, et comment ce découpage se lit sur la matrice ? Si
$E=F\oplus G$ et si un endomorphisme $f$ envoie $F$ dans $F$ et $G$ dans $G$, sa matrice dans une base adaptée est formée de deux
blocs indépendants : on étudie $f$ morceau par morceau. Ce chapitre fournit le vocabulaire de ce découpage (sommes directes,
hyperplans, sous-espaces stables) et un premier invariant qui ne dépend pas de la base, la trace.

\ouverture{%
\begin{itemize}[leftmargin=1.1em,itemsep=1pt]
\item Maîtriser l'algèbre linéaire de première année : bases, dimension, rang, matrices, changement de base (§\ref{ss:rappels}).
\item Démontrer qu'une somme est directe ; obtenir une base adaptée et la dimension d'une somme directe.
\item Reconnaître un projecteur ou une symétrie et en écrire la matrice dans une base adaptée.
\item Passer d'un hyperplan ou d'un sous-espace à ses équations, et réciproquement.
\item Utiliser la trace et ses propriétés ; lire des sous-espaces stables sur une matrice par blocs.
\end{itemize}}{%
Algèbre linéaire de première année : espaces vectoriels, familles libres et génératrices, bases, dimension, applications
linéaires et matrices (rappelés au §\ref{ss:rappels}).}{Deux séances}

Dans tout le chapitre, $\K=\R$ ou $\C$ et $E$ est un $\K$-espace vectoriel ; $\LE$ est l'ensemble des endomorphismes de $E$.
Quand $E$ est de dimension finie, on note $n=\dim E$.

% ---------------------------------------------------------------------
\subsection{Rappels de première année}\label{ss:rappels}
% Numérotation propre aux rappels (R.1, R.2, ...) ; la numérotation du chapitre reprend à 1.1 au §1.2.
\renewcommand{\thedefinition}{R.\arabic{definition}}
\makeatletter\renewcommand{\thetcb@cnt@methode}{R.\arabic{tcb@cnt@methode}}\makeatother
\renewcommand{\theexo}{R.\arabic{exo}}

\begin{commencer}
Cette partie rassemble tout le programme d'algèbre linéaire de première année (PTSI) : calcul matriciel et systèmes, espaces
vectoriels et dimension, applications linéaires, matrices et changements de base. Ces résultats sont utilisés dans tout le
parcours d'algèbre ; ils sont énoncés sans démonstration (numérotés R.1, R.2\dots). Les sommes directes, les supplémentaires,
les projecteurs et les symétries, également vus en première année, sont repris au §\ref{ss:sommes} pour un nombre quelconque de
sous-espaces. Le déterminant fait l'objet du chapitre~2.
\end{commencer}

\subsubsection*{Calcul matriciel et systèmes linéaires}

\begin{definition}[opérations sur les matrices]
$\mathcal M_{n,p}(\K)$ est l'ensemble des matrices à $n$ lignes et $p$ colonnes ; $E_{i,j}$ désigne la matrice dont le seul
coefficient non nul est un $1$ en position $(i,j)$ (\emph{matrice élémentaire}). Pour $A\in\mathcal M_{n,p}(\K)$ et
$B\in\mathcal M_{p,q}(\K)$, le produit $AB\in\mathcal M_{n,q}(\K)$ a pour coefficients
\[(AB)_{i,j}=\sum_{k=1}^pa_{i,k}b_{k,j}.\]
La \emph{transposée} $A^T\in\mathcal M_{p,n}(\K)$ est définie par $(A^T)_{i,j}=a_{j,i}$.
\end{definition}

\begin{proposition}[règles de calcul]\label{r:calcul}
\begin{enumerate}
\item Le produit est associatif et bilinéaire ; $I_nA=AI_p=A$. Il n'est pas commutatif, et $AB=0$ n'entraîne pas $A=0$ ou $B=0$.
\item $(A^T)^T=A$, $(\lambda A+\mu B)^T=\lambda A^T+\mu B^T$ et $(AB)^T=B^TA^T$.
\item \emph{Formule du binôme.} Si $A,B\in\Mn{\K}$ \emph{commutent} ($AB=BA$), pour tout $m\in\N$ :
$\displaystyle(A+B)^m=\sum_{k=0}^m\binom mkA^kB^{m-k}$.
\item Les matrices diagonales, triangulaires supérieures, triangulaires inférieures forment des sous-espaces stables par
produit ; les coefficients diagonaux d'un produit de deux matrices triangulaires supérieures sont les produits des coefficients
diagonaux. $A$ est \emph{symétrique} si $A^T=A$, \emph{antisymétrique} si $A^T=-A$.
\end{enumerate}
\end{proposition}

\begin{proposition}[matrices inversibles]\label{r:inverse}
$A\in\Mn{\K}$ est \emph{inversible} s'il existe $B\in\Mn{\K}$ telle que $AB=BA=I_n$ ; on note $B=A^{-1}$ et $\GLn{\K}$ l'ensemble
des matrices inversibles (le \emph{groupe linéaire}).
\begin{enumerate}
\item Si $A,B\in\GLn{\K}$ : $AB\in\GLn{\K}$ et $(AB)^{-1}=B^{-1}A^{-1}$ ; $A^T\in\GLn{\K}$ et $(A^T)^{-1}=(A^{-1})^T$.
\item Il suffit d'une seule des deux égalités : si $AB=I_n$, alors $A$ et $B$ sont inversibles et $B=A^{-1}$.
\item Une matrice triangulaire est inversible si et seulement si ses coefficients diagonaux sont tous non nuls.
\item $\begin{pmatrix}a&b\\c&d\end{pmatrix}$ est inversible si et seulement si $ad-bc\neq0$, et son inverse est
$\dfrac1{ad-bc}\begin{pmatrix}d&-b\\-c&a\end{pmatrix}$.
\end{enumerate}
\end{proposition}

\begin{proposition}[opérations élémentaires et systèmes]\label{r:systemes}
\begin{enumerate}
\item Les opérations élémentaires sur les lignes ($L_i\leftrightarrow L_j$ ; $L_i\leftarrow\lambda L_i$ avec $\lambda\neq0$ ;
$L_i\leftarrow L_i+\lambda L_j$ avec $i\neq j$) reviennent à multiplier \emph{à gauche} par une matrice inversible ; sur les
colonnes, à multiplier \emph{à droite}.
\item Un système linéaire s'écrit $AX=B$, avec $A\in\mathcal M_{n,p}(\K)$, $X\in\mathcal M_{p,1}(\K)$ l'inconnue et
$B\in\mathcal M_{n,1}(\K)$. Des opérations élémentaires sur les lignes du système le transforment en un système
\emph{équivalent} (algorithme du pivot de Gauss).
\item Si le système est \emph{compatible} (il a au moins une solution $X_0$), ses solutions sont les $X_0+Y$, où $Y$ décrit
l'ensemble des solutions du système homogène $AY=0$.
\end{enumerate}
\end{proposition}

\begin{methode}{Calculer l'inverse d'une matrice}\label{m:inverse}
\begin{itemize}
\item \emph{Par un système :} résoudre $AX=Y$ d'inconnue $X$ pour $Y$ quelconque ; si l'on trouve $X=BY$ pour tout $Y$, alors
$A$ est inversible et $A^{-1}=B$.
\item \emph{Par le pivot :} appliquer à $I_n$ les opérations sur les lignes qui transforment $A$ en $I_n$ ; on obtient $A^{-1}$.
Si l'on fait apparaître une ligne nulle, $A$ n'est pas inversible.
\item \emph{Par un polynôme annulateur :} si $A^2+aA+bI_n=0$ avec $b\neq0$, alors $A\big(-\tfrac1b(A+aI_n)\big)=I_n$.
\end{itemize}
\end{methode}

\begin{exemple}
Soit $A=\begin{pmatrix}1&1&0\\0&1&1\\0&0&1\end{pmatrix}$. Le système $AX=Y$ s'écrit $x_1+x_2=y_1$, $x_2+x_3=y_2$, $x_3=y_3$ ; en
remontant : $x_3=y_3$, $x_2=y_2-y_3$, $x_1=y_1-y_2+y_3$. Donc $A$ est inversible et
$A^{-1}=\begin{pmatrix}1&-1&1\\0&1&-1\\0&0&1\end{pmatrix}$.
\end{exemple}

\subsubsection*{Espaces vectoriels, familles de vecteurs}

\begin{definition}[espaces et sous-espaces vectoriels]
Les espaces vectoriels de référence sont $\K^n$, $\mathcal M_{n,p}(\K)$, $\K[X]$ et $\K_n[X]$ (polynômes de degré au plus $n$),
l'espace des fonctions d'un ensemble $\Omega$ dans $\K$ (ou dans un espace vectoriel), et l'espace $\K^\N$ des suites. Une
\emph{combinaison linéaire} de $x_1,\dots,x_p$ est un vecteur $\lambda_1x_1+\dots+\lambda_px_p$. Une partie $F$ de $E$ est un
\emph{sous-espace vectoriel} si $0\in F$ et si $\lambda x+\mu y\in F$ pour tous $x,y\in F$ et $\lambda,\mu\in\K$.
\end{definition}

\begin{proposition}[sous-espace engendré]\label{r:vect}
\begin{enumerate}
\item Une intersection de sous-espaces vectoriels est un sous-espace vectoriel (une réunion, en général, non).
\item $\vect(x_1,\dots,x_p)$, l'ensemble des combinaisons linéaires de $x_1,\dots,x_p$, est un sous-espace vectoriel : c'est
le plus petit qui contient $x_1,\dots,x_p$ (tout sous-espace qui contient les $x_i$ le contient).
\end{enumerate}
\end{proposition}

\begin{definition}[familles libres, génératrices, bases]
Une famille $(x_1,\dots,x_p)$ de $E$ est
\begin{itemize}
\item \emph{génératrice} si $\vect(x_1,\dots,x_p)=E$ ;
\item \emph{libre} si $\lambda_1x_1+\dots+\lambda_px_p=0$ entraîne $\lambda_1=\dots=\lambda_p=0$ ; \emph{liée} sinon ;
\item une \emph{base} si elle est libre et génératrice : tout $x\in E$ s'écrit alors de manière unique
$x=\lambda_1x_1+\dots+\lambda_px_p$, et $(\lambda_1,\dots,\lambda_p)$ sont les \emph{coordonnées} de $x$.
\end{itemize}
Bases canoniques : $(e_1,\dots,e_n)$ pour $\K^n$, $(E_{i,j})$ pour $\mathcal M_{n,p}(\K)$, $(1,X,\dots,X^n)$ pour $\K_n[X]$.
\end{definition}

\begin{proposition}[familles libres]\label{r:libres}
\begin{enumerate}
\item Une famille est liée si et seulement si l'un de ses vecteurs est combinaison linéaire des autres.
\item Si $(x_1,\dots,x_p)$ est libre et $x_{p+1}\notin\vect(x_1,\dots,x_p)$, alors $(x_1,\dots,x_{p+1})$ est libre.
\item Une famille de polynômes non nuls de degrés deux à deux distincts est libre ; si $\deg P_k=k$ pour $0\leq k\leq n$,
$(P_0,\dots,P_n)$ est une base de $\K_n[X]$ (base \emph{échelonnée en degré}).
\end{enumerate}
\end{proposition}

\begin{methode}{Montrer qu'une famille est libre, ou est une base}\label{m:libre}
Écrire une combinaison linéaire nulle et en déduire, par un système, que tous les coefficients sont nuls. En dimension $n$, une
famille de $n$ vecteurs est une base dès qu'elle est libre (ou dès qu'elle est génératrice) : une seule vérification suffit
(proposition~\ref{r:base}). Dans $\K^n$, on peut aussi calculer le rang de la matrice des vecteurs en colonnes, ou un
déterminant (chapitre~2).
\end{methode}

\subsubsection*{Dimension finie}

\begin{theoreme}[bases et dimension]\label{r:dim}
On suppose que $E$ possède une famille génératrice finie ($E$ est \emph{de dimension finie}).
\begin{enumerate}
\item \emph{Base extraite.} De toute famille génératrice de $E$, on peut extraire une base de $E$.
\item \emph{Base incomplète.} Toute famille libre de $E$ peut être complétée en une base de $E$ (à l'aide de vecteurs d'une
famille génératrice donnée).
\item Dans un espace engendré par $n$ vecteurs, toute famille de $n+1$ vecteurs est liée. Par suite, toutes les bases de $E$ ont
le même nombre d'éléments, la \emph{dimension} de $E$.
\end{enumerate}
$\dim\K^n=n$, \ $\dim\K_n[X]=n+1$, \ $\dim\mathcal M_{n,p}(\K)=np$, \ $\dim\mathcal L(E,F)=\dim E\times\dim F$.
\end{theoreme}

\begin{proposition}[familles en dimension $n$]\label{r:base}
Si $\dim E=n$ : une famille libre a au plus $n$ vecteurs, une famille génératrice au moins $n$, et pour une famille de
\emph{$n$ vecteurs},
\[\text{base}\iff\text{libre}\iff\text{génératrice}.\]
Le \emph{rang} d'une famille est $\rg(x_1,\dots,x_p)=\dim\vect(x_1,\dots,x_p)$ ; il vaut au plus $p$, avec égalité si et
seulement si la famille est libre.
\end{proposition}

\begin{proposition}[sous-espaces et dimension]\label{r:sev}
Soient $F$ et $G$ deux sous-espaces de $E$, de dimension finie.
\begin{enumerate}
\item $\dim F\leq\dim E$, avec égalité si et seulement si $F=E$.
\item \emph{Formule de Grassmann :} $\dim(F+G)=\dim F+\dim G-\dim(F\cap G)$.
\item Tout sous-espace $F$ de $E$ possède au moins un supplémentaire (en général, une infinité).
\item $E=F\oplus G$ si et seulement si $F\cap G=\{0\}$ et $\dim F+\dim G=\dim E$.
\end{enumerate}
\end{proposition}

\begin{exemple}[espaces de solutions]
Les ensembles suivants sont des espaces vectoriels de dimension finie, dont on connaît la dimension :
\begin{itemize}
\item les solutions sur un intervalle de $y'+a(t)y=0$ ($a$ continue) : dimension $1$ ;
\item les solutions de $y''+ay'+by=0$ ($a,b$ constants) : dimension $2$ (chapitre~4) ;
\item les suites vérifiant $u_{n+2}=au_{n+1}+bu_n$ : dimension $2$, car $u\mapsto(u_0,u_1)$ est un isomorphisme sur $\K^2$.
\end{itemize}
\end{exemple}

\subsubsection*{Applications linéaires}

\begin{definition}[applications linéaires]
$u:E\to F$ est \emph{linéaire} si $u(\lambda x+\mu y)=\lambda u(x)+\mu u(y)$ ; leur ensemble $\mathcal L(E,F)$ est un espace
vectoriel, et la composée de deux applications linéaires est linéaire ($v\circ u$ est bilinéaire en $(u,v)$). Un
\emph{isomorphisme} est une application linéaire bijective (sa réciproque est linéaire) ; un \emph{endomorphisme} est un élément
de $\LE=\mathcal L(E,E)$, un \emph{automorphisme} un endomorphisme bijectif, et $\mathrm{GL}(E)$ est l'ensemble des
automorphismes. On note $\Id_E$ l'identité, $\lambda\Id_E$ les homothéties, et $u^k=u\circ\dots\circ u$ ($u^0=\Id_E$ ; $k\in\Z$
si $u\in\mathrm{GL}(E)$).
\end{definition}

\begin{proposition}[noyau, image, rang]\label{p:rang}
Soit $u:E\to F$ linéaire.
\begin{enumerate}
\item L'image directe et l'image réciproque d'un sous-espace par $u$ sont des sous-espaces. En particulier,
$\Ker u=\{x\in E,\ u(x)=0\}$ et $\Im u=u(E)$ sont des sous-espaces vectoriels ; $u$ est injective si et seulement si
$\Ker u=\{0\}$, surjective si et seulement si $\Im u=F$.
\item Si $(e_1,\dots,e_p)$ engendre $E$, alors $\Im u=\vect\big(u(e_1),\dots,u(e_p)\big)$.
\item \emph{Théorème du rang.} Si $E$ est de dimension finie, $\dim\Ker u+\rg u=\dim E$, où $\rg u=\dim\Im u$. Plus précisément,
tout supplémentaire de $\Ker u$ dans $E$ est isomorphe à $\Im u$ (par la restriction de $u$).
\item $\rg(v\circ u)\leq\min(\rg u,\rg v)$ ; composer par un isomorphisme (à gauche ou à droite) ne change pas le rang.
\item Si $\dim E=\dim F$ (finie), en particulier pour un endomorphisme : $u$ injective $\iff$ $u$ surjective $\iff$ $u$ bijective.
Un endomorphisme inversible à gauche (ou à droite) est donc inversible.
\end{enumerate}
\end{proposition}

\begin{proposition}[une application linéaire est déterminée par une base]\label{r:base-image}
Soit $(e_1,\dots,e_n)$ une base de $E$.
\begin{enumerate}
\item Pour toute famille $(f_1,\dots,f_n)$ de $F$, il existe une unique $u\in\mathcal L(E,F)$ telle que $u(e_i)=f_i$ pour tout $i$.
\item $u$ est injective (resp.\ surjective, bijective) si et seulement si $\big(u(e_1),\dots,u(e_n)\big)$ est libre
(resp.\ génératrice de $F$, une base de $F$).
\item Deux espaces de dimension finie sont isomorphes si et seulement s'ils ont la même dimension.
\item Si $E=F_1\oplus F_2$, deux applications linéaires définies sur $F_1$ et sur $F_2$ se prolongent de manière unique en une
application linéaire définie sur $E$.
\end{enumerate}
\end{proposition}

\begin{proposition}[équations linéaires]\label{r:equations}
Soient $u\in\mathcal L(E,F)$ et $a\in F$. L'ensemble des solutions de l'équation $u(x)=a$ est soit vide, soit de la forme
\[x_0+\Ker u=\{x_0+h,\ h\in\Ker u\},\]
où $x_0$ est une solution particulière. Exemples : un système $AX=B$ (proposition~\ref{r:systemes}) ; une équation
différentielle linéaire, $S_E=y_p+S_H$ (chapitre~4) ; une suite arithmético-géométrique $u_{n+1}=au_n+b$ ($a\neq1$), somme de la
suite constante $\frac b{1-a}$ et d'une suite géométrique de raison $a$.
\end{proposition}

\subsubsection*{Matrices et applications linéaires}

\begin{proposition}[matrice d'une application linéaire]\label{r:matrice}
Soient $\mathcal B=(e_1,\dots,e_p)$ une base de $E$ et $\mathcal C$ une base de $F$ ($\dim F=n$). La matrice
$\Mat_{\mathcal B,\mathcal C}(u)\in\mathcal M_{n,p}(\K)$ a pour $j$-ième colonne les coordonnées de $u(e_j)$ dans $\mathcal C$.
\begin{enumerate}
\item Si $X$ est la colonne des coordonnées de $x$ dans $\mathcal B$ et $Y$ celle de $u(x)$ dans $\mathcal C$, alors $Y=AX$.
\item $u\mapsto\Mat_{\mathcal B,\mathcal C}(u)$ est un isomorphisme de $\mathcal L(E,F)$ sur $\mathcal M_{n,p}(\K)$ ; la matrice d'une
composée est le produit des matrices : $\Mat(v\circ u)=\Mat(v)\Mat(u)$ (dans des bases compatibles).
\item $u$ est un isomorphisme si et seulement si sa matrice (carrée) est inversible, et $\Mat(u^{-1})=\Mat(u)^{-1}$.
\end{enumerate}
\end{proposition}

\begin{proposition}[noyau, image et rang d'une matrice]\label{r:rang-matrice}
On identifie $\mathcal M_{p,1}(\K)$ à $\K^p$. L'application $X\mapsto AX$ est l'application linéaire \emph{canoniquement
associée} à $A\in\mathcal M_{n,p}(\K)$ ; son noyau, son image et son rang sont ceux de $A$.
\begin{enumerate}
\item $\Im A$ est engendrée par les colonnes de $A$ ; les lignes de $A$ fournissent un système d'équations de $\Ker A$.
\item Le rang de $A$ est le rang de la famille de ses colonnes, et le rang de toute application linéaire représentée par $A$.
\item Les opérations élémentaires sur les colonnes conservent l'image et le rang ; sur les lignes, le noyau et le rang.
\item $\rg(A^T)=\rg A$ (admis).
\item Pour $A\in\Mn{\K}$ : $A$ inversible $\iff\Ker A=\{0\}\iff\rg A=n\iff$ les colonnes de $A$ forment une base de $\K^n$.
\end{enumerate}
\end{proposition}

\begin{methode}{Calculer un rang, un noyau, une image}\label{m:rang}
\begin{itemize}
\item \emph{Rang :} par des opérations élémentaires, se ramener à une matrice échelonnée ; le rang est le nombre de pivots
(lignes non nulles).
\item \emph{Noyau :} résoudre le système $AX=0$ ; sa dimension est $p-\rg A$ (théorème du rang), ce qui permet de contrôler le
résultat.
\item \emph{Image :} elle est engendrée par les colonnes ; on en extrait une base de $\rg A$ colonnes non colinéaires
(par exemple celles des pivots).
\end{itemize}
\end{methode}

\begin{exemple}
$A=\begin{pmatrix}1&2&1\\2&4&0\\1&2&-1\end{pmatrix}$. La deuxième colonne vaut deux fois la première, et les colonnes $C_1$, $C_3$
ne sont pas colinéaires : $\rg A=2$ et $\Im A=\vect\big((1,2,1),(1,0,-1)\big)$. Par le théorème du rang, $\dim\Ker A=1$ ; comme
$2C_1-C_2=0$, $\Ker A=\vect\big((2,-1,0)\big)$.
\end{exemple}

\begin{proposition}[changement de base]\label{p:chgt}
Soient $\mathcal B$ et $\mathcal B'$ deux bases de $E$ ($\dim E=n$). La matrice de passage $P=P_{\mathcal B\to\mathcal B'}$ a pour
colonnes les coordonnées des vecteurs de $\mathcal B'$ dans $\mathcal B$ ; elle est inversible, d'inverse $P_{\mathcal B'\to\mathcal B}$.
\begin{enumerate}
\item Si $X$ et $X'$ sont les coordonnées d'un même vecteur dans $\mathcal B$ et $\mathcal B'$, alors $X=PX'$.
\item Pour $f\in\LE$ : $\Mat_{\mathcal B'}(f)=P^{-1}\,\Mat_{\mathcal B}(f)\,P$. Le rang de $f$ est celui de n'importe laquelle de ses
matrices.
\item Plus généralement, pour $u\in\mathcal L(E,F)$, si $Q$ est la matrice de passage entre deux bases de $F$ :
$A'=Q^{-1}AP$.
\end{enumerate}
\end{proposition}

\begin{definition}
Deux matrices $A,B\in\Mn{\K}$ sont \emph{semblables} s'il existe $P\in\GLn{\K}$ telle que $B=P^{-1}AP$ : elles représentent le
même endomorphisme dans deux bases. Chercher une base où la matrice d'un endomorphisme est simple, c'est chercher une matrice
simple semblable à une matrice donnée : c'est l'objet du chapitre~3.
\end{definition}

\begin{proposition}[systèmes linéaires et rang]\label{r:cramer}
Soit $A\in\mathcal M_{n,p}(\K)$.
\begin{enumerate}
\item L'ensemble des solutions de $AX=0$ est $\Ker A$, de dimension $p-\rg A$.
\item Le système $AX=B$ est compatible si et seulement si $B\in\Im A$.
\item Si $A\in\GLn{\K}$ (système \emph{de Cramer}), $AX=B$ a une unique solution, $X=A^{-1}B$.
\end{enumerate}
\end{proposition}

\begin{reflexe}
Tout ce qui ne dépend que de l'endomorphisme, et pas de la base, est commun à deux matrices semblables : le rang, la trace
(§\ref{ss:trace}), le déterminant (chapitre~2), le polynôme caractéristique et les valeurs propres (chapitre~3). Si deux matrices
diffèrent par l'un de ces invariants, elles ne sont pas semblables.
\end{reflexe}

\begin{exo}{al-r-base}{1}{base et coordonnées}
Dans $\R^3$, on pose $u_1=(1,1,0)$, $u_2=(0,1,1)$, $u_3=(1,0,1)$.
\begin{enumerate}
\item Montrer que $\mathcal B'=(u_1,u_2,u_3)$ est une base de $\R^3$.
\item Déterminer les coordonnées de $v=(2,3,1)$ dans $\mathcal B'$.
\end{enumerate}
\end{exo}
\begin{exo}{al-r-rang}{2}{noyau, image et théorème du rang}
Soit $f:\R^3\to\R^3$, $f(x,y,z)=(x+y-z,\ x+2y,\ y+z)$.
\begin{enumerate}
\item Écrire la matrice $A$ de $f$ dans la base canonique.
\item Déterminer une base de $\Ker f$ et une base de $\Im f$, et vérifier le théorème du rang.
\item $f$ est-elle injective ? surjective ? Donner une équation de $\Im f$.
\end{enumerate}
\end{exo}
\begin{exo}{al-r-chgt}{2}{inverse et changement de base}
Soit $f$ l'endomorphisme de $\R^2$ de matrice $A=\begin{pmatrix}3&-2\\1&0\end{pmatrix}$ dans la base canonique, et
$u_1=(1,1)$, $u_2=(2,1)$.
\begin{enumerate}
\item Montrer que $\mathcal B'=(u_1,u_2)$ est une base de $\R^2$ et calculer la matrice de passage $P$ et son inverse.
\item Calculer $f(u_1)$ et $f(u_2)$, puis la matrice $D$ de $f$ dans $\mathcal B'$. Vérifier que $D=P^{-1}AP$.
\item En déduire $A^k$ pour $k\in\N$.
\end{enumerate}
\end{exo}

% fin des rappels : la numérotation du chapitre reprend
\renewcommand{\thedefinition}{\thesection.\arabic{definition}}\setcounter{definition}{0}
\makeatletter\renewcommand{\thetcb@cnt@methode}{\thesection.\arabic{tcb@cnt@methode}}\setcounter{tcb@cnt@methode}{0}\makeatother
\renewcommand{\theexo}{\thesection.\arabic{exo}}\setcounter{exo}{0}

% ---------------------------------------------------------------------
\subsection{Produit et somme de sous-espaces}\label{ss:sommes}

\subsubsection*{Produit d'espaces vectoriels}

\begin{proposition}[produit]\label{p:produit}
Si $E_1,\dots,E_p$ sont des $\K$-espaces vectoriels, le produit $E_1\times\dots\times E_p$, muni des opérations coordonnée par
coordonnée, est un $\K$-espace vectoriel. Si chaque $E_i$ est de dimension finie,
\[\dim(E_1\times\dots\times E_p)=\dim E_1+\dots+\dim E_p.\]
\end{proposition}

\begin{proof}
Pour $p=2$ : si $(a_1,\dots,a_n)$ est une base de $E_1$ et $(b_1,\dots,b_m)$ une base de $E_2$, la famille
$\big((a_1,0),\dots,(a_n,0),(0,b_1),\dots,(0,b_m)\big)$ est une base de $E_1\times E_2$ : tout couple $(x,y)$ s'écrit
$(x,0)+(0,y)$, et une combinaison nulle donne, coordonnée par coordonnée, deux combinaisons nulles. Puis récurrence sur $p$.
\end{proof}

\begin{exemple}
$\dim(\R^2\times\R_2[X])=2+3=5$. L'ensemble des solutions $(x,y)\in\R\times\R$ d'un système est une partie de $\R^2=\R\times\R$ :
c'est le premier exemple de produit.
\end{exemple}

\subsubsection*{Somme de plusieurs sous-espaces, somme directe}

\begin{definition}
Soient $F_1,\dots,F_p$ des sous-espaces vectoriels de $E$. Leur \emph{somme} est
$F_1+\dots+F_p=\{x_1+\dots+x_p,\ x_i\in F_i\}$ : c'est un sous-espace vectoriel, le plus petit qui contient tous les $F_i$.
La somme est \emph{directe}, et l'on note $F_1\oplus\dots\oplus F_p$, si tout vecteur de la somme s'écrit de manière
\emph{unique} $x_1+\dots+x_p$ avec $x_i\in F_i$.
\end{definition}

\begin{proposition}[caractérisation]\label{p:directe}
La somme $F_1+\dots+F_p$ est directe si et seulement si
\[x_1+\dots+x_p=0\text{ avec }x_i\in F_i\quad\Longrightarrow\quad x_1=\dots=x_p=0.\]
Pour deux sous-espaces : $F+G$ est directe si et seulement si $F\cap G=\{0\}$.
\end{proposition}

\begin{proof}
Si la somme est directe, $0=0+\dots+0$ est l'unique décomposition de $0$. Réciproquement, si $x_1+\dots+x_p=y_1+\dots+y_p$,
alors $(x_1-y_1)+\dots+(x_p-y_p)=0$ avec $x_i-y_i\in F_i$, donc $x_i=y_i$. Pour $p=2$ : si $u\in F\cap G$, $u+(-u)=0$ avec
$u\in F$, $-u\in G$, donc $u=0$ ; réciproquement, si $x+y=0$ avec $x\in F$, $y\in G$, alors $x=-y\in F\cap G=\{0\}$.
\end{proof}

\begin{attention}
Pour $p\geq3$, il ne suffit pas que les intersections deux à deux soient nulles : dans $\R^2$, les trois droites
$\vect(1,0)$, $\vect(0,1)$, $\vect(1,1)$ se coupent deux à deux en $\{0\}$, mais $(1,0)+(0,1)-(1,1)=0$ : la somme n'est pas directe.
La seule caractérisation au programme est celle de la proposition~\ref{p:directe}.
\end{attention}

\begin{theoreme}[base adaptée, dimension]\label{t:adaptee}
On suppose $E$ de dimension finie. Si la somme $F_1+\dots+F_p$ est directe, la concaténation de bases des $F_i$ est une base
de $F_1\oplus\dots\oplus F_p$ (base \emph{adaptée}), et
\[\dim(F_1\oplus\dots\oplus F_p)=\dim F_1+\dots+\dim F_p.\]
En particulier, $E=F_1\oplus\dots\oplus F_p$ si et seulement si la concaténation de bases des $F_i$ est une base de $E$.
\end{theoreme}

\begin{proof}
La famille concaténée engendre la somme. Elle est libre : une combinaison linéaire nulle s'écrit $x_1+\dots+x_p=0$, où $x_i$
regroupe les termes de la base de $F_i$ ; la somme étant directe, chaque $x_i$ est nul, puis chaque coefficient (base de $F_i$).
\end{proof}

Pour deux sous-espaces, la première année a fourni la \emph{formule de Grassmann}, $\dim(F+G)=\dim F+\dim G-\dim(F\cap G)$
(proposition~\ref{r:sev}), dont on déduit la caractérisation pratique des supplémentaires.

\begin{enpratique}
Pour montrer que $E=F\oplus G$ en dimension finie, il suffit de deux des trois propriétés : $F\cap G=\{0\}$, $F+G=E$,
$\dim F+\dim G=\dim E$. Le plus souvent : intersection nulle et somme des dimensions.
\end{enpratique}

\subsubsection*{Projecteurs et symétries}

Une somme directe $E=F\oplus G$ fournit deux endomorphismes : pour $x=x_F+x_G$ (décomposition unique), la \emph{projection}
$p(x)=x_F$ sur $F$ parallèlement à $G$, et la \emph{symétrie} $s(x)=x_F-x_G$ par rapport à $F$ parallèlement à $G$. Ils sont
linéaires, et $s=2p-\Id_E$.

\begin{proposition}[caractérisation, rappel]\label{p:projsym}
Soit $f\in\LE$.
\begin{enumerate}
\item $f$ est un projecteur si et seulement si $f\circ f=f$ ; c'est alors la projection sur $\Im f$ parallèlement à $\Ker f$, et
$E=\Im f\oplus\Ker f$.
\item $f$ est une symétrie si et seulement si $f\circ f=\Id_E$ ; c'est alors la symétrie par rapport à $\Ker(f-\Id_E)$ parallèlement
à $\Ker(f+\Id_E)$, et $E=\Ker(f-\Id_E)\oplus\Ker(f+\Id_E)$.
\end{enumerate}
\end{proposition}

\begin{enpratique}
Dans une base adaptée à ces sommes directes, avec $r=\rg p$ et $r=\dim\Ker(s-\Id_E)$ :
\[\Mat(p)=\begin{pmatrix}I_r&0\\0&0\end{pmatrix},\qquad\Mat(s)=\begin{pmatrix}I_r&0\\0&-I_{n-r}\end{pmatrix}.\]
Pour calculer $p(x)$, on écrit $x=x_F+x_G$ : on cherche la composante sur $G$ par la condition « $x-x_G\in F$ ».
\end{enpratique}

\begin{exo}{red-somme}{1}{sommes directes}
\begin{enumerate}
\item Dans $\R^3$, on pose $F=\{(x,y,z),\ x+y+z=0\}$ et $G=\vect\big((1,1,1)\big)$. Montrer que $\R^3=F\oplus G$ et décomposer
$v=(x,y,z)$ selon cette somme.
\item Dans $\R^3$, les droites $D_1=\vect(1,0,0)$, $D_2=\vect(0,1,0)$, $D_3=\vect(1,1,0)$ sont-elles en somme directe ?
\end{enumerate}
\end{exo}
\begin{exo}{al-proj}{2}{projecteur et symétrie dans $\R^3$}
Dans $\R^3$, on considère le plan $P$ d'équation $2x-y+z=0$ et la droite $D$ d'équations $y=0$, $x=z$.
\begin{enumerate}
\item Donner une base de $P$ et de $D$, puis montrer que $\R^3=P\oplus D$.
\item Déterminer la matrice, dans la base canonique, de la projection $p$ sur $P$ parallèlement à $D$, puis celle de la symétrie $s$
par rapport à $P$ parallèlement à $D$. Vérifier que $p\circ p=p$ et $s\circ s=\Id$.
\end{enumerate}
\end{exo}

% ---------------------------------------------------------------------
\subsection{Hyperplans et systèmes linéaires}\label{ss:hyperplans}

\begin{definition}
Un \emph{hyperplan} de $E$ est un sous-espace qui admet une droite comme supplémentaire. En dimension finie $n\geq1$, c'est un
sous-espace de dimension $n-1$ (théorème~\ref{t:adaptee}).
\end{definition}

\begin{proposition}\label{p:hyperplan}
On suppose $E$ de dimension finie $n\geq1$, muni d'une base $\mathcal B$.
\begin{enumerate}
\item Les hyperplans de $E$ sont exactement les noyaux des formes linéaires non nulles $\varphi:E\to\K$. En coordonnées dans
$\mathcal B$, un hyperplan a une équation $a_1x_1+\dots+a_nx_n=0$ avec $(a_1,\dots,a_n)\neq(0,\dots,0)$, unique à un facteur non
nul près.
\item Si $H=\Ker\varphi$ et $u\notin H$, alors $E=H\oplus\vect(u)$.
\item L'intersection de $p$ hyperplans est de dimension au moins $n-p$. Réciproquement, tout sous-espace de dimension $n-p$
est l'intersection de $p$ hyperplans : il est défini par un système de $p$ équations.
\end{enumerate}
\end{proposition}

\begin{proof}
1. Si $\varphi\neq0$, son image est $\K$ (seul sous-espace non nul de $\K$), donc $\dim\Ker\varphi=n-1$ par le théorème du rang.
Réciproquement, si $\dim H=n-1$, on complète une base $(e_1,\dots,e_{n-1})$ de $H$ en une base de $E$ par $e_n$, et la forme
« $n$-ième coordonnée » a pour noyau $H$. En coordonnées, $\varphi(x)=\sum a_ix_i$ avec $a_i=\varphi(e_i)$.
2. $u\notin H$ donne $H\cap\vect(u)=\{0\}$, et $\dim H+1=n$.
3. Par récurrence : si $G$ est un sous-espace et $H=\Ker\varphi$, la restriction de $\varphi$ à $G$ a un rang au plus $1$, donc
$\dim(G\cap H)\geq\dim G-1$. Réciproquement, on complète une base $(e_1,\dots,e_{n-p})$ du sous-espace en une base de $E$ :
il est défini par la nullité des $p$ dernières coordonnées.
\end{proof}

\begin{enpratique}
L'ensemble des solutions d'un système linéaire homogène de $p$ équations à $n$ inconnues est un sous-espace de $\K^n$ de dimension
$n-r$, où $r$ est le rang du système (théorème du rang appliqué à $X\mapsto AX$) : chaque équation « indépendante » fait perdre
une dimension. L'ensemble des solutions d'un système \emph{avec second membre} est vide, ou de la forme $X_0+\Ker A$ : un sous-espace
affine dirigé par l'ensemble des solutions du système homogène. Dans $\R^3$, un plan vectoriel a une équation, une droite vectorielle deux.
\end{enpratique}

\begin{methode}{Passer d'une base à des équations, et inversement}\label{m:equations}
\begin{itemize}
\item Des équations à une base : résoudre le système (inconnues principales en fonction des inconnues libres).
\item D'une base $(v_1,\dots,v_k)$ de $F$ à des équations : écrire que $x\in F\iff$ le système $x=\lambda_1v_1+\dots+\lambda_kv_k$
d'inconnues $\lambda_i$ est compatible ; les conditions de compatibilité sont les $n-k$ équations de $F$.
\end{itemize}
\end{methode}

\begin{exemple}
$F=\vect\big((1,2,3,4)\big)$ dans $\R^4$ : $(x,y,z,t)\in F\iff y=2x,\ z=3x,\ t=4x$. Trois équations indépendantes : $F$ est
l'intersection des trois hyperplans $y-2x=0$, $z-3x=0$, $t-4x=0$ (et $1=4-3$).
\end{exemple}

\begin{exo}{red-hyp}{2}{hyperplans de $\R^4$}
Dans $\R^4$, on considère $H_1$ : $x+y+z+t=0$ et $H_2$ : $x-y+z-t=0$, et $u=(1,1,1,1)$.
\begin{enumerate}
\item Donner la dimension et une base de $H_1\cap H_2$.
\item La droite $\vect(u)$ est-elle un supplémentaire de $H_1$ ? de $H_2$ ?
\end{enumerate}
\end{exo}

% ---------------------------------------------------------------------
\subsection{Trace}\label{ss:trace}

\begin{definition}
La \emph{trace} de $A=(a_{i,j})\in\Mn{\K}$ est la somme de ses coefficients diagonaux : $\tr(A)=\sum_{i=1}^na_{i,i}$.
\end{definition}

\begin{proposition}\label{p:trace}
\begin{enumerate}
\item La trace est une forme linéaire non nulle sur $\Mn{\K}$ : son noyau, les matrices de trace nulle, est un hyperplan de
$\Mn{\K}$, de dimension $n^2-1$.
\item Pour $A\in\mathcal M_{n,p}(\K)$ et $B\in\mathcal M_{p,n}(\K)$ : $\tr(AB)=\tr(BA)$.
\item Deux matrices semblables ont la même trace. On définit donc la trace d'un endomorphisme $f$ en dimension finie :
$\tr(f)=\tr\big(\Mat_{\mathcal B}(f)\big)$, indépendante de la base $\mathcal B$.
\end{enumerate}
\end{proposition}

\begin{proof}
2. $\tr(AB)=\sum_{i=1}^n\sum_{k=1}^pa_{i,k}b_{k,i}=\sum_{k=1}^p\sum_{i=1}^nb_{k,i}a_{i,k}=\tr(BA)$.
3. $\tr(P^{-1}AP)=\tr\big((P^{-1}A)P\big)=\tr\big(P(P^{-1}A)\big)=\tr(A)$, par le point 2.
\end{proof}

\begin{attention}
$\tr(AB)\neq\tr(A)\tr(B)$ en général, et $\tr(ABC)=\tr(CAB)$ (permutation circulaire) mais $\tr(ABC)\neq\tr(BAC)$ en général.
\end{attention}

\begin{exemple}[trace d'un projecteur, d'une symétrie]
Si $p$ est un projecteur de rang $r$, dans une base adaptée à $E=\Im p\oplus\Ker p$, sa matrice est
$\mathrm{diag}(I_r,0)$ : $\tr(p)=r=\rg(p)$. Pour une symétrie $s$ : $\tr(s)=\dim\Ker(s-\Id)-\dim\Ker(s+\Id)$.
\end{exemple}

\begin{exo}{red-trace}{2}{trace}
\begin{enumerate}
\item Montrer qu'il n'existe pas de matrices $A,B\in\Mn{\K}$ telles que $AB-BA=I_n$.
\item Soit $A\in\Mn{\K}$ telle que $A^2=A$. Montrer que $\rg(A)=\tr(A)$.
\end{enumerate}
\end{exo}

% ---------------------------------------------------------------------
\subsection{Sous-espaces stables}\label{ss:stables}

\begin{definition}
Un sous-espace $F$ de $E$ est \emph{stable} par $f\in\LE$ si $f(F)\subset F$. L'application $f_F:F\to F$, $x\mapsto f(x)$, est alors un
endomorphisme de $F$, l'\emph{endomorphisme induit} par $f$ sur $F$.
\end{definition}

\begin{proposition}\label{p:stable}
\begin{enumerate}
\item Si $F=\vect(e_1,\dots,e_p)$, $F$ est stable par $f$ si et seulement si $f(e_j)\in F$ pour tout $j$ : on vérifie sur une
famille génératrice.
\item $\Ker f$ et $\Im f$ sont stables par $f$. Plus généralement, si $f\circ g=g\circ f$, alors $\Ker g$ et $\Im g$ sont stables par $f$.
\end{enumerate}
\end{proposition}

\begin{proof}
1. Si $x=\sum\lambda_je_j\in F$, $f(x)=\sum\lambda_jf(e_j)\in F$. 2. Si $g(x)=0$, $g\big(f(x)\big)=f\big(g(x)\big)=0$ : $f(x)\in\Ker g$.
Si $y=g(x)$, $f(y)=g\big(f(x)\big)\in\Im g$. Pour $g=f$, on obtient la stabilité de $\Ker f$ et $\Im f$.
\end{proof}

\begin{theoreme}[traduction matricielle]\label{t:blocs}
On suppose $E$ de dimension finie. Soit $F$ un sous-espace de dimension $p$ et $\mathcal B=(e_1,\dots,e_n)$ une base adaptée à $F$
($F=\vect(e_1,\dots,e_p)$).
\begin{enumerate}
\item $F$ est stable par $f$ si et seulement si $\Mat_{\mathcal B}(f)$ est triangulaire par blocs :
$\begin{pmatrix}A&B\\0&C\end{pmatrix}$, avec $A\in\mathcal M_p(\K)$ ; $A$ est alors la matrice de l'endomorphisme induit $f_F$.
\item Si $E=F\oplus G$ et si $\mathcal B$ est adaptée à cette somme, $F$ et $G$ sont stables par $f$ si et seulement si
$\Mat_{\mathcal B}(f)$ est diagonale par blocs : $\begin{pmatrix}A&0\\0&C\end{pmatrix}$.
\end{enumerate}
Dans les deux cas, $\tr f=\tr A+\tr C$.
\end{theoreme}

\begin{proof}
La $j$-ième colonne contient les coordonnées de $f(e_j)$. Pour $j\leq p$, $f(e_j)\in F$ si et seulement si ses coordonnées sur
$e_{p+1},\dots,e_n$ sont nulles : c'est le bloc $0$. Pour 2, on applique 1 à $F$ et à $G$. La trace se lit sur la diagonale.
\end{proof}

\begin{enpratique}
Lire une matrice par blocs, c'est lire des sous-espaces stables, et réciproquement. Un zéro en bas à gauche dit : « les $p$
premiers vecteurs de base engendrent un sous-espace stable ». Les projecteurs et les symétries (§\ref{ss:sommes}) en sont les
premiers exemples : $\Im p$ et $\Ker p$ sont stables par $p$, et sa matrice est diagonale par blocs. Le chapitre~2 ajoutera
$\det f=\det A\times\det C$.
\end{enpratique}

\begin{proposition}[droites stables]\label{p:droite}
Soit $u\neq0$. La droite $\vect(u)$ est stable par $f$ si et seulement s'il existe $\lambda\in\K$ tel que $f(u)=\lambda u$.
\end{proposition}

\begin{proof}
$\vect(u)$ est stable si et seulement si $f(u)\in\vect(u)$, c'est-à-dire $f(u)=\lambda u$ pour un scalaire $\lambda$.
\end{proof}

Cette dernière proposition ouvre le chapitre \emph{Réduction} : chercher les droites stables, c'est chercher les vecteurs que $f$
se contente de dilater (les \emph{vecteurs propres}). S'il y en a assez pour former une base, la matrice de $f$ dans cette base
est diagonale : c'est le cas extrême d'une matrice diagonale par blocs, avec des blocs de taille $1$.

\begin{exo}{red-stable}{2}{lire des blocs}
Soit $f$ l'endomorphisme de $\R^3$ de matrice $A=\begin{pmatrix}3&-1&1\\0&2&0\\1&-1&3\end{pmatrix}$ dans la base canonique
$(e_1,e_2,e_3)$.
\begin{enumerate}
\item Montrer que le plan $F$ d'équation $y=0$ et la droite $G=\vect(u)$, $u=(1,1,0)$, sont stables par $f$.
\item Montrer que $\R^3=F\oplus G$ et écrire la matrice de $f$ dans la base $(e_1,e_3,u)$. Vérifier sur cette matrice que $\tr A=8$.
\end{enumerate}
\end{exo}
\begin{exo}{al-crochet}{3}{l'endomorphisme $M\mapsto AM-MA$}
Soit $A=\begin{pmatrix}1&1\\1&-1\end{pmatrix}$ et $\varphi:\mathcal M_2(\R)\to\mathcal M_2(\R)$, $M\mapsto AM-MA$.
\begin{enumerate}
\item Montrer que $\varphi$ est un endomorphisme et que son image est contenue dans l'hyperplan des matrices de trace nulle.
\item Écrire la matrice de $\varphi$ dans la base canonique $(E_{1,1},E_{1,2},E_{2,1},E_{2,2})$. Déterminer $\Ker\varphi$ et $\Im\varphi$.
\item Montrer que $\Ker\varphi$ et $\Im\varphi$ sont stables par $\varphi$. Sont-ils supplémentaires ?
\end{enumerate}
\end{exo}

% ---------------------------------------------------------------------
\subsection{Quelle méthode ? Synthèse}\label{ss:synthese-al}

\begin{center}
\renewcommand{\arraystretch}{1.35}
\small
\begin{tabular}{p{0.38\linewidth}p{0.55\linewidth}}
\hline
\textbf{Question} & \textbf{Méthode}\\\hline
$\dim\Ker u$, $\rg u$ ? & une base de l'image (vecteurs $u(e_j)$), puis le théorème du rang\\
$F+G$ directe ? & $F\cap G=\{0\}$ (proposition~\ref{p:directe}) ; pour $p\geq3$ sous-espaces : unicité de la décomposition de $0$\\
$E=F\oplus G$ ? & $F\cap G=\{0\}$ et $\dim F+\dim G=\dim E$ (théorème~\ref{t:adaptee})\\
$f$ projecteur, symétrie ? & $f\circ f=f$, $f\circ f=\Id$ ; base adaptée à $\Im f\oplus\Ker f$ ou $\Ker(f-\Id)\oplus\Ker(f+\Id)$\\
calculer une projection & décomposer $x=x_F+x_G$ : la condition $x-x_G\in F$ donne $x_G$\\
équations d'un sous-espace & condition de compatibilité de $x=\sum\lambda_iv_i$ (méthode~\ref{m:equations}) ; $n-\dim F$ équations\\
trace & linéaire, $\tr(AB)=\tr(BA)$, invariante par similitude ; $\tr p=\rg p$\\
$F$ stable par $f$ ? & vérifier $f(e_j)\in F$ sur une base de $F$ ; matrice triangulaire par blocs (théorème~\ref{t:blocs})\\\hline
\end{tabular}
\end{center}

\begin{attention}\textbf{Pièges fréquents.}
\begin{itemize}[leftmargin=1.1em]
\item Somme directe de trois sous-espaces : les intersections deux à deux ne suffisent pas.
\item $E=F\oplus G$ demande aussi la somme (ou les dimensions), pas seulement $F\cap G=\{0\}$.
\item La matrice de passage a pour colonnes les \emph{nouveaux} vecteurs exprimés dans l'\emph{ancienne} base, et
$\Mat_{\mathcal B'}(f)=P^{-1}AP$, pas $PAP^{-1}$.
\item $\Ker f$ et $\Im f$ ne sont pas supplémentaires en général (penser à $f\neq0$ avec $f\circ f=0$).
\item $\tr(AB)\neq\tr A\,\tr B$ ; $\tr(ABC)\neq\tr(BAC)$ en général.
\end{itemize}
\end{attention}

\begin{essentiel}
\begin{itemize}[leftmargin=1.1em,itemsep=2pt]
\item Théorème du rang ; $\Mat_{\mathcal B'}(f)=P^{-1}\Mat_{\mathcal B}(f)P$ ; matrices semblables : même rang, même trace.
\item Somme directe $\iff$ unicité de la décomposition de $0$ ; base adaptée, $\dim\bigoplus F_i=\sum\dim F_i$.
\item Projecteur $\iff p\circ p=p$, $E=\Im p\oplus\Ker p$ ; symétrie $\iff s\circ s=\Id$.
\item Hyperplan $=$ noyau d'une forme linéaire non nulle $=$ sous-espace de dimension $n-1$ ; $p$ équations $\Rightarrow$ dimension $\geq n-p$.
\item $\tr(AB)=\tr(BA)$ ; $F$ stable $\iff$ matrice triangulaire par blocs ; droite stable $\iff$ $f(u)=\lambda u$.
\end{itemize}
\end{essentiel}

% =====================================================================
%                              CORRIGÉS
% =====================================================================
\clearpage
\renewcommand{\ongletcouleur}{Ocre}
\section{Corrigés}\label{sec:corriges-al}

\begin{corrige}{al-r-base}
\q{1.} Si $au_1+bu_2+cu_3=0$ : $a+c=0$, $a+b=0$, $b+c=0$. Donc $c=-a$, $b=-a$, puis $b+c=-2a=0$ : $a=b=c=0$. La famille est
libre, et c'est une famille de $3=\dim\R^3$ vecteurs : c'est une base (proposition~\ref{r:base}).

\q{2.} On cherche $(a,b,c)$ avec $a+c=2$, $a+b=3$, $b+c=1$. En additionnant : $2(a+b+c)=6$, soit $a+b+c=3$ ; d'où $b=3-2=1$,
$c=3-3=0$, $a=3-1=2$. Les coordonnées de $v$ dans $\mathcal B'$ sont $(2,1,0)$ : $v=2u_1+u_2$.
\end{corrige}

\begin{corrige}{al-r-rang}
\q{1.} Les colonnes sont les images des vecteurs de la base canonique :
$A=\begin{pmatrix}1&1&-1\\1&2&0\\0&1&1\end{pmatrix}$.

\q{2.} $\Ker f$ : $x+y-z=0$, $x+2y=0$, $y+z=0$. Les deux dernières équations donnent $x=-2y$, $z=-y$, et la première est
alors vérifiée ($-2y+y+y=0$). Donc $\Ker f=\vect\big((-2,1,-1)\big)$, de dimension $1$.\\
$\Im f$ est engendrée par les colonnes $C_1=(1,1,0)$, $C_2=(1,2,1)$, $C_3=(-1,0,1)$ ; le vecteur du noyau donne la relation
$-2C_1+C_2-C_3=0$, donc $C_3\in\vect(C_1,C_2)$, et $C_1$, $C_2$ ne sont pas colinéaires : $(C_1,C_2)$ est une base de $\Im f$.
Théorème du rang : $\dim\Ker f+\rg f=1+2=3=\dim\R^3$.

\q{3.} $\Ker f\neq\{0\}$ : $f$ n'est pas injective ; $\rg f=2<3$ : elle n'est pas surjective (ce qui était prévisible : pour un
endomorphisme en dimension finie, les deux vont ensemble). Une équation $\alpha x+\beta y+\gamma z=0$ de $\Im f$ doit être vérifiée
par $C_1$ et $C_2$ : $\alpha+\beta=0$, $\alpha+2\beta+\gamma=0$ ; avec $\alpha=1$ : $\beta=-1$, $\gamma=1$. Donc
$\Im f$ : $x-y+z=0$ (et $C_3$ la vérifie bien).
\end{corrige}

\begin{corrige}{al-r-chgt}
\q{1.} $u_1$ et $u_2$ ne sont pas colinéaires : c'est une famille libre de $2$ vecteurs de $\R^2$, donc une base.
$P=\begin{pmatrix}1&2\\1&1\end{pmatrix}$, $1\times1-2\times1=-1\neq0$, et
$P^{-1}=\dfrac1{-1}\begin{pmatrix}1&-2\\-1&1\end{pmatrix}=\begin{pmatrix}-1&2\\1&-1\end{pmatrix}$.

\q{2.} $f(u_1)=(3-2,\,1)=(1,1)=u_1$ et $f(u_2)=(6-2,\,2)=(4,2)=2u_2$. Donc $D=\begin{pmatrix}1&0\\0&2\end{pmatrix}$. Vérification :
$AP=\begin{pmatrix}1&4\\1&2\end{pmatrix}$ et $P^{-1}(AP)=\begin{pmatrix}-1+2&-4+4\\1-1&4-2\end{pmatrix}=D$.

\q{3.} $A=PDP^{-1}$, donc $A^k=PD^kP^{-1}$ (les $P^{-1}P$ intermédiaires se simplifient), avec $D^k=\begin{pmatrix}1&0\\0&2^k\end{pmatrix}$ :
\[A^k=\begin{pmatrix}1&2^{k+1}\\1&2^k\end{pmatrix}\begin{pmatrix}-1&2\\1&-1\end{pmatrix}=\begin{pmatrix}2^{k+1}-1&2-2^{k+1}\\2^k-1&2-2^k\end{pmatrix}.\]
Contrôle : $k=1$ redonne $A$, $k=0$ redonne $I_2$.
\end{corrige}

\begin{corrige}{red-somme}
\q{1.} Si $v\in F\cap G$, $v=(m,m,m)$ et $3m=0$ : $F\cap G=\{0\}$. De plus $\dim F=2$ (un hyperplan de $\R^3$) et $\dim G=1$ : la
somme est directe et de dimension $3$, donc $\R^3=F\oplus G$. On cherche $v=w+m(1,1,1)$ avec $w\in F$ : la somme des coordonnées
de $w$ est nulle, d'où $x+y+z=3m$. Donc
\[v=\underbrace{\big(x-m,\,y-m,\,z-m\big)}_{\in F}+\underbrace{m(1,1,1)}_{\in G},\qquad m=\frac{x+y+z}3.\]
\q{2.} Les intersections deux à deux sont nulles, mais $(1,0,0)+(0,1,0)-(1,1,0)=0$ est une décomposition non triviale de $0$ :
la somme n'est pas directe (proposition~\ref{p:directe}). D'ailleurs $D_1+D_2+D_3$ est le plan $z=0$, de dimension $2<3$.
\end{corrige}

\begin{corrige}{al-proj}
\choix{on décompose un vecteur quelconque selon $P\oplus D$ : la condition « $x-x_D\in P$ » donne directement la composante sur $D$.}
\q{1.} $P$ : $y=2x+z$, de base $\big((1,2,0),(0,1,1)\big)$. $D$ : $(x,0,x)$, de base $w=(1,0,1)$. Comme $2\cdot1-0+1=3\neq0$,
$w\notin P$ : $P\cap D=\{0\}$, et $\dim P+\dim D=3$, donc $\R^3=P\oplus D$ (c'est aussi la proposition~\ref{p:hyperplan}, point 2).
\q{2.} Soit $v=(x,y,z)$. On cherche $t$ tel que $v-tw\in P$ : $2(x-t)-y+(z-t)=0$, soit $t=\frac{2x-y+z}3$. Alors
$p(v)=v-tw$ et $s(v)=2p(v)-v=v-2tw$. Dans la base canonique :
\[\Mat(p)=\frac13\begin{pmatrix}1&1&-1\\0&3&0\\-2&1&2\end{pmatrix},\qquad\Mat(s)=\frac13\begin{pmatrix}-1&2&-2\\0&3&0\\-4&2&1\end{pmatrix}.\]
Le calcul direct donne $\Mat(p)^2=\Mat(p)$ et $\Mat(s)^2=I_3$. Contrôles : $\tr p=2=\rg p=\dim P$, et $\tr s=2-1=1$.
\retenir{pour une projection, on ne cherche que la composante « parallèle » ; la symétrie s'en déduit par $s=2p-\Id$.}
\end{corrige}

\begin{corrige}{red-hyp}
\q{1.} En ajoutant et en retranchant les équations : $x+z=0$ et $y+t=0$. Donc $H_1\cap H_2=\{(x,y,-x,-y)\}$, de base
$\big((1,0,-1,0),(0,1,0,-1)\big)$ : dimension $2=4-2$.
\q{2.} $u\notin H_1$ ($1+1+1+1=4\neq0$) : $\R^4=H_1\oplus\vect(u)$ (proposition~\ref{p:hyperplan}). En revanche $u\in H_2$
($1-1+1-1=0$) : $\vect(u)\subset H_2$, ce n'est pas un supplémentaire.
\end{corrige}

\begin{corrige}{red-trace}
\q{1.} Par linéarité et $\tr(AB)=\tr(BA)$ : $\tr(AB-BA)=0$, alors que $\tr(I_n)=n\neq0$.
\q{2.} L'endomorphisme $p$ de $\K^n$ associé à $A$ est un projecteur ; dans une base adaptée à $\K^n=\Im p\oplus\Ker p$, sa matrice
est $\mathrm{diag}(I_r,0)$ avec $r=\rg(p)$. Donc $A$ est semblable à cette matrice, et $\tr A=r=\rg A$.
\retenir{la trace ne voit pas l'ordre d'un produit de deux matrices : $\tr(AB-BA)=0$.}
\end{corrige}

\begin{corrige}{red-stable}
\q{1.} $f(x,0,z)=(3x+z,\,0,\,x+3z)$ a une deuxième coordonnée nulle : $F$ est stable. $f(u)=A(1,1,0)^T=(2,2,0)=2u$ : $G$ est stable.
\q{2.} $u\notin F$ (sa deuxième coordonnée vaut $1$) : $\R^3=F\oplus G$. Dans la base $(e_1,e_3,u)$ : $f(e_1)=(3,0,1)=3e_1+e_3$,
$f(e_3)=(1,0,3)=e_1+3e_3$, $f(u)=2u$, donc
\[\Mat(f)=\begin{pmatrix}3&1&0\\1&3&0\\0&0&2\end{pmatrix},\qquad\tr A=3+3+2=8.\]
\retenir{deux sous-espaces stables supplémentaires $\Rightarrow$ matrice diagonale par blocs.}
\end{corrige}

\begin{corrige}{al-crochet}
\choix{on calcule l'image de chaque vecteur de la base canonique ; la trace donne le contrôle de l'image.}
\q{1.} $\varphi(\lambda M+N)=A(\lambda M+N)-(\lambda M+N)A=\lambda\varphi(M)+\varphi(N)$, et $\varphi(M)\in\mathcal M_2(\R)$. Par
$\tr(AM)=\tr(MA)$, $\tr\varphi(M)=0$ : $\Im\varphi$ est contenue dans l'hyperplan $\Ker\tr$, de dimension $3$.
\q{2.} $\varphi(E_{1,1})=\begin{pmatrix}0&-1\\1&0\end{pmatrix}$, $\varphi(E_{1,2})=\begin{pmatrix}-1&2\\0&1\end{pmatrix}$,
$\varphi(E_{2,1})=\begin{pmatrix}1&0\\-2&-1\end{pmatrix}$, $\varphi(E_{2,2})=\begin{pmatrix}0&1\\-1&0\end{pmatrix}$, d'où
\[\Mat(\varphi)=\begin{pmatrix}0&-1&1&0\\-1&2&0&1\\1&0&-2&-1\\0&1&-1&0\end{pmatrix}.\]
Les colonnes $C_1$ et $C_2$ ne sont pas proportionnelles, $C_4=-C_1$ et $C_3=-2C_1-C_2$ : le rang vaut $2$, donc
$\dim\Ker\varphi=2$ (théorème du rang). Or $I_2$ et $A$ commutent avec $A$, et ne sont pas proportionnelles : $\Ker\varphi=\vect(I_2,A)$.
Et $\Im\varphi=\vect\big(\varphi(E_{1,1}),\varphi(E_{1,2})\big)$, un plan contenu dans l'hyperplan des matrices de trace nulle.
\q{3.} $\Ker\varphi$ et $\Im\varphi$ sont stables par $\varphi$ (proposition~\ref{p:stable}). Une matrice de $\Ker\varphi\cap\Im\varphi$
s'écrit $aI_2+bA$ avec une trace nulle : $2a=0$ ($\tr A=0$), donc $bA\in\Im\varphi$. Or $A=\begin{pmatrix}1&1\\1&-1\end{pmatrix}$ n'est
pas dans $\Im\varphi$ : si $A=\lambda\varphi(E_{1,1})+\mu\varphi(E_{1,2})=\begin{pmatrix}-\mu&-\lambda+2\mu\\\lambda&\mu\end{pmatrix}$,
le coefficient $(2,1)$ donne $\lambda=1$, le coefficient $(1,1)$ donne $\mu=-1$, et le coefficient $(1,2)$ vaudrait $-3\neq1$. Donc $b=0$ : l'intersection est
nulle, et $\dim\Ker\varphi+\dim\Im\varphi=4$ : ils sont supplémentaires.
\retenir{$M\mapsto AM-MA$ : son noyau contient toujours $\vect(I_n,A)$, son image est dans l'hyperplan des matrices de trace nulle.}
\end{corrige}

\clearpage
\ongletsfalse
\section*{Index des méthodes}\label{sec:index}
\addcontentsline{toc}{section}{Index des méthodes}
\markboth{Index des méthodes}{}
{\small\listemethodes}
\end{document}
